\subsection{Critical checks on the Berezinian and flatness signs}

These calculations audit faithful mathematical errors in the source. They are
elementary critical annotations, not a claim to new research and not a development
of a class C or D problem.

\paragraph{The reciprocal Berezinian on P000673.}\oursComputed{}
For an even block-diagonal matrix \(J=\operatorname{diag}(a,d)\) of size
\(1|1\), with nonzero even scalar entries, the usual Berezinian is
\(\operatorname{Ber}J=a/d\). The first expression on P000673 gives \(d/a\),
whereas the lower expression on the same page gives \(a/d\).
Taking \(a=2,d=3\) gives \(3/2\) versus \(2/3\), so the two expressions
cannot be simultaneous definitions in the stated convention.
The usual normalization is explicitly printed in
\citet[\S1]{bernsteinleites1977integral} and in the earlier density formula
\citep[formula~(3)]{berezinleites1975supermanifolds}.
The correct general formula is
\[
 \operatorname{Ber}\begin{pmatrix}A&B\\C&D\end{pmatrix}
 =\frac{\det(A-BD^{-1}C)}{\det D}
 =\frac{\det A}{\det(D-CA^{-1}B)}.
\]
This test isolates a reciprocal error without depending on conventions for
left versus right odd derivatives.

\paragraph{Flatness on P000684 and in the printed Manin theorem.}\ours{}
Use the notebook's convention
\(\nabla_1\mathbf1=P\mathbf1\) and
\(\nabla_2\mathbf1=L\mathbf1\).
Because both connections commute with \(\nabla\), the coefficientwise
Leibniz rule gives
\[
 [\nabla_1,\nabla_2]\mathbf1
 = (d_1L+LP-d_2P-PL)\mathbf1.
\]
The cyclic filtered-module theorem makes the action on \(\mathbf1\)
faithful for differential operators, hence flatness implies
\[
 d_1L-d_2P=[P,L].
\]
Both the difference sign and the order of the commutator follow from the
declared convention. An unspecified \(\pm\) in front of a commutator cannot
repair a sum of the two derivatives.

For a concrete counterexample to that sum, take
\(\mathcal B=\mathbb Q[t,z]\), \(d=0\),
\(\mathcal M=\mathcal B[D]\), and
\[
 \nabla Q=DQ,\qquad
 \nabla_1Q=\partial_tQ+Qz,\qquad
 \nabla_2Q=\partial_zQ+Qt.
\]
The standard order filtration satisfies the module axioms; the three
connections commute. Their values on \(1\) are \(P=z\) and \(L=t\).
Thus \(d_1L-d_2P=1-1=0=[P,L]\), whereas
\(d_1L+d_2P=2\ne0\). The accompanying exact symbolic check verifies this
on a general test polynomial in \(t,z,D\).

The error is also present in the inspected printed English statement
\citep[Chapter~III, Theorem~2.6, p.~83]{manin1978algebraic}.
That statement uses the opposite names,
\(\nabla_1\mathbf1=L\mathbf1\), \(\nabla_2\mathbf1=P\mathbf1\),
and prints \(d_1P+d_2L=[P,L]\). Its proof requires instead
\(d_1P-d_2L=[L,P]\). This identifies a shared mathematical sign error,
not a failure to transcribe the notebook.

\paragraph{Independent fourth-order check.}\oursComputed{}
For P000681, the companion's separate \texttt{critical\_checks.py} computes
the commutator by two methods. With \(L=D^4+vD^2+wD+z\),
\(P=D^2+AD+v/2\), and \(A_x=0\), its coefficients are
\[
\begin{aligned}
[L,P]_{D^2}&=2v_{xx}-Av_x-2w_x,\\
[L,P]_{D}&=2v_{xxx}+vv_x-w_{xx}-Aw_x-2z_x,\\
[L,P]_{1}&=\tfrac12v_{xxxx}+\tfrac12vv_{xx}+\tfrac12wv_x-z_{xx}-Az_x.
\end{aligned}
\]
The notebook's first two relative coefficients are therefore wrong, while
the third equation agrees. Manin's later equation has the corrected relative
coefficients, with its own time and commutator conventions
\citep[Chapter~III, \S6.3, equation~(17), p.~96]{manin1978algebraic}.
